\subsection{The actual abelian quotients at boundary points}
\label{sec:actual-point-tangent-degrees}

The flat determinant construction applies to the abelian quotient at
an actual boundary point. This supplies the signed degree comparison
used in the point-height argument. The integral image of the tangent
lattice need not be saturated; its index has a favourable sign and is
retained below.

Keep the original host $E$, with $f=[E^+:\mathbb Q]$, the prescribed
quadratic subfield $\mathbb Q(\sqrt{-d})$, and the hypotheses that
$-d$ is the original odd fundamental discriminant and
$E/E^+$ is unramified at every finite prime. Retain the original
primitive element $w=(1+\sqrt{-d})/2$, with $w+\bar w=1$. Put
\[
 P_E=\gamma+\log(2\pi)-\log\operatorname{rd}(E)
             -\frac{L'(1,\chi_{-d})}{L(1,\chi_{-d})}.
\]
Let an algebraic point of the original compactification have toric
$\mathcal O_E$-rank $r$ and signature $(p,q)$, with
$0\le r\le\min(p,q)$. Its abelian quotient $A$ has signature
$(p-r,q-r)$. After a finite extension $L$ containing the coefficient
embeddings and admitting semistable reduction, let $\bar V$ be the
Lie lattice of the N\'eron identity component of $A$, with the
prescribed metric $\|v\|^2=\pi H(v,v)$. At every conjugate pair of
embeddings, write $W_\tau,W'_\tau$ for its saturated integral
eigenlattices, with inherited metrics. The chosen half has rank
$m-r$, where $m$ is either $p$ or $q$; the other half has rank
$p+q-m-r$. All degrees below are divided by $[L:\mathbb Q]$.

\begin{theorem}[Conjugate degrees at the actual points]
\label{thm:actual-point-conjugate-degrees}
For these lattices,
\begin{equation}\label{eq:actual-point-conjugate-degrees}
 \widehat{\deg}_{\rm n}W_\tau-
 \widehat{\deg}_{\rm n}W'_\tau
                    =\frac{2m-p-q}{2}P_E.
\end{equation}
This is an identity on the actual integral lattices, including the
places ramified in $E$ and the places above two. It is unchanged by
further finite extension. The extension used to realize the point
does not introduce a field-size term in this normalized identity.
\end{theorem}

\begin{proof}
\emph{The quotient and its polarization.}
The original primitive line and polarization restrict to the
semiabelian fibre. Its toric character module is a projective
$\mathcal O_E$-module of rank $r$, so every coefficient embedding
occurs $r$ times in its Lie algebra. Exactness of the Lie sequence
in characteristic zero gives the stated quotient signature. The
principal degeneration formula
\[
                 1=[X:\phi(Y)]^2\deg\lambda_A
\]
makes $\lambda_A$ principal. The original primitive identity gives
the rigidified line
$(1,\lambda_A\iota(w))^*\mathcal P_A$, with
$w+\bar w=1$. These are precisely the quotient and line
identifications in the arithmetic-point uniformization theorem
\cite{current-construction-record}; no independent polarization or
translate of its line is chosen.

\emph{Realization in the appropriate moduli space.}
If $p+q-2r=0$, both eigenlattices are zero and the assertion follows.
Otherwise fix one complex embedding and take the integral Betti
lattice of this principally polarized quotient, with its coefficient
action and perfect alternating form. Betti--\'etale comparison gives
the required complete Tate genus. The moduli construction is rerun
for this quotient's own integral Betti and Tate datum. Choose an
allowed odd half-level,
and pass to a finite field extension over which its finite torsion
labels and the required roots are defined. The generic PEL moduli
theorem then represents this very abelian variety, primitive line and
both entire finite labels. The component-field and global-model
constructions apply to its signature, including either definite
signature \cite{current-construction-record}. Their bounds are not
being used to bound the field of definition of the point. Passing
to the auxiliary cover may enlarge $L$ again; this is permitted
because the conclusion is invariant under normalized semistable
base change, with no bound on that extension required.

The global model is proper. Thus the resulting point extends over
every valuation ring of its number field. Its semiabelian family
has the same generic fibre as the prescribed semistable model of
$A$. The semistable N\'eron comparison, with the zero section fixed,
identifies their identity components and their integral Lie modules
\cite[Theorem 4.4 and Corollary 4.5]{conrad-semistable}.
The rigidified restriction argument in the theta construction
identifies the original lines and their canonical metrics. Hence
the integral eigenlattices and their Hermitian norms are exactly
the ones in the statement. This argument also applies after
ramified field extension; it does not compare only generic fibres.

\emph{Evaluation of the same determinant.}
Apply the actual flat-degree theorem to this component and to the
signature $(a,b)=(m-r,p+q-m-r)$. Its corrected flat degree is
\[
 \frac{a-b}{2}P_E+b(\log d-\log\operatorname{rd}(E)).
\]
On integral eigenframes, the prescribed finite logarithmic norms
add $b(\log|\mathfrak D_E|-\log|d|)$. Their normalized sum is
$b(\log d-\log\operatorname{rd}(E))$. Degree is the negative
sum of logarithmic norms, so this term cancels the displayed finite
correction. The exact determinant comparison
\eqref{geo-exact-conjugate-degree} therefore gives
$(a-b)P_E/2$ for the actual eigenlattice degree difference.
Since $a-b=2m-p-q$, this is
\eqref{eq:actual-point-conjugate-degrees}.

Eigenlattices are kernels of the finitely many integral coefficient
action equations. Flat base change preserves these kernels.
Semistable base change preserves the Lie lattice, and normalized
Hermitian degree is unchanged. This proves the field-extension
assertion, with every original finite correction retained.
\end{proof}

Let $V_\tau\subset W_\tau$ be the image of the original integral
tangent lattice in the selected quotient eigenspace. The quotient
map extends to the N\'eron model by its mapping property, so the
inclusion is integral. It has full generic rank. Set
\[
 g=f(p+q-2r),\qquad R=f(m-r),\qquad
 T=\sum_\tau\widehat{\deg}_{\rm n}V_\tau,\qquad
 b=\frac{f(2m-p-q)}2P_E.
\]
The coefficient summands are orthogonal for the prescribed Hodge
metric. Let $\delta\ge0$ be the normalized logarithmic index of
$\bigoplus_\tau(W_\tau\oplus W'_\tau)$ in $V$, and let
$\nu\ge0$ be the sum of the normalized logarithmic indices
$[W_\tau:V_\tau]$.

\begin{corollary}[The selected degrees for the actual family]
\label{cor:actual-selected-point-degrees}
For the original lattices,
\begin{equation}\label{eq:actual-selected-point-degrees}
 T\le12R\log\bigl(2+g+(b-\delta-2\nu)_+\bigr).
\end{equation}
If $B\ge\max\{128,4f(p+q)^2\}$ and
\[
 \Lambda_B=\log B+\log\operatorname{rd}(E)
                  +\left|\frac{L'(1,\chi_{-d})}{L(1,\chi_{-d})}\right|,
\]
then
\begin{equation}\label{eq:actual-selected-point-budget}
                       T\le36f(m-r)\Lambda_B.
\end{equation}
Neither statement assumes a bound on either integral index.
\end{corollary}

\begin{proof}
Write $d_V=\widehat{\deg}_{\rm n}V$. Orthogonality and the two
integral inclusions give
\[
                        d_V=2T-b+\delta+2\nu.
\]
The classical theta construction supplies the tangent-slope estimate
with constant $12$ on this exact principally polarized semistable
lattice and its metric. Applied to the selected integral images,
it gives $T\le12R\log(2+g+(-d_V)_+)$. If $T>0$, the displayed
identity gives $(-d_V)_+\le(b-\delta-2\nu)_+$; if $T\le0$,
\eqref{eq:actual-selected-point-degrees} is immediate. The case
$g=0$ has $T=R=0$ and needs no slope theorem.

The previously established period estimate is
$|P_E|\le C_f\Lambda_B$, with
$C_f=1+\gamma+|\log(2\pi)|<4$. Moreover
$g\le G=f(p+q)\le B\le e^{\Lambda_B}$ and
$|f(2m-p-q)|\le g$. Thus $b\le C_fG\Lambda_B$.
The selected-degree calculation in
Proposition~\ref{geo-selected-tangent-sum} gives
$T\le12\{1+\log(3+C_f)\}R\Lambda_B<36R\Lambda_B$.
The weak inequality in \eqref{eq:actual-selected-point-budget}
includes all zero-rank cases.
\end{proof}

This removes the conjugate-degree and selected-abelian-degree
hypotheses from the actual point-height application. The corrected fibre comparison is supplied in
Section~\ref{sec:geometry-hodge-specialization}. The regularized
point-height identity and final mean-volume estimate are separate
obligations; they do not follow from a tangent-degree bound alone.
